%!TEX program = xelatex
\documentclass[cn,hazy,blue,14pt,screen]{elegantnote}
\title{2020级物理第一次班会}
\author{冷轩}

%\version{2.30}
\date{\zhtoday}



\begin{document}
\maketitle

\centerline{
  \includegraphics[width=0.4\textwidth]{nbu1.png}
}

\newpage
\tableofcontents   \label{content}
\phantomsection%防止目录中页码和pdf链接中的页码不一致，而超链失效。



\newpage
\section{班委选举}

{\centering \Large 班委选举}


\newpage
\section{班主任的定位}
自己对大学班主任角色的定位，就两点：
\begin{itemize}
\item 站在学生的角度上，和班上学生一条阵线；
\begin{itemize}
\item 不整虚的，讲实在话；
\item ... ...
\end{itemize}

\item 辅助班上学生顺利毕业
\begin{itemize}
\item 辅助走得高的，走的更高远；辅助走得慢的，跟上；
\item 帮助引进提升的机会，如比赛和科研；
\item ... ...
\end{itemize}

\item 教课和班主任基本是用爱发电，珍惜关心你们的老师~~
\end{itemize}



\newpage
\section{大学生的几点重要的思维转变}
\subsection{成年人意识}
{\bf 大学生基本是准成年人和成年人了。}

\begin{itemize}
\item  成年人的第一个要点：为自己的行为负责，为自己的前途命运负责。

\item  进而引出一个重要的意识：主动。

\item  主动：让自己忙碌起来，时间用在刀刃上、主动抓住机会。
\end{itemize}


{\bf 外部环境发生变化：父母不在身边、老师不再贴身督促。}
\begin{itemize}
\item  进而引出另一个重要的意识：自觉、自律
\end{itemize}

 
 {\bf 做人做事的态度（很多时候机会只有一次）}
 \begin{itemize}
\item 思考未来，思考自己的位置；
 
\item 第一印象很关键；

\item 守时是基本；

\item 打开格局。
\end{itemize}
 
 {\bf 一开始就要有毕业愿景}（要有大方向，不清晰没关系，及时明确）
 
 
 
 
 \newpage
 \subsubsection{问“问题”与听取建议、意见}
\paragraph{听取建议、意见}
夫子说过“三人行必有我师焉”。无论是前辈、同学、同事、晚辈沟通，都有听取意见的时候。听取他人意见不是听取他人对自己的观点的附和，往往是与自己观点不同甚至相背的。这时候就让人很难受。但是难受是让自己进步的前兆。如果非常坚持自己的观点，那就给他人观点一个机会，先试用一下他人观点，比较一下优缺点。最终结果往往不是自己或他人观点的胜出，而是不同观点的融合。所以要经常反思：我给过别人建议、意见机会了吗？

\paragraph{RTFM} is an initialism for the expression "read the fucking manual" – typically used to reply to a question that could have been easily answered from the product manual or documentation.\footnote{\url{https://en.wikipedia.org/wiki/RTFM}}

大家的时间和精力都是有限的，请不要问些“简单问题”浪费别人的生命。每个人都应该具备自主学习能力，不要张口就要“答案”，这是不利于自己成长的。问别人“问题”前请反问自己两点：
\begin{itemize}
\item 网上搜索过了吗？
\item 自己思考过了吗？
\item 流程：首先尝试自己解决，再和同事讨论，再问上级。
\item 记录解决问题过程的心得体会，然后自己解决问题的能力与速度会不断提升，进而关注到别人未解决的问题，此时你已经站到人类知识的边界了，为人类拓荒了。
\end{itemize}

 
 

\newpage
\subsection{课程学习}
{ \bf 四字核心：详略得当}
\begin{itemize}
\item 发达国家名校的课程学习：一般课上松，考试难

\item 除申请国外研究生需要刷分数绩点外，毕业就业和考研一般只认学位，几乎没有关心你本科课程分数的

\item 人的精力是有限的，以自己未来发展方向为中心，合理调整各个课程花费的时间，水课划水，自己核心课认真学习

\item 全面扎实的基本概念最重要，并不需要特别深入，也不会像高考一样变着花子考你。在你入职或升学后，自然而然会要你深入

\item 注重积累与沉淀，形成自己的东西（记笔记和积累程序是个好方法，记笔记推荐用\LaTeX ）
\end{itemize}


{\bf 英语要过六级，兼顾学习生活英语（参考我写的文档）}


{\bf 作为物理班，要会一门编程语言，推荐Python，无论是工作中还是科研中都非常强大}


\newpage
\paragraph{素描学习法}
这种方法是从网络视频，普通物理-台湾国立交通大学-李威仪，那学来的。其理论依据应该可以用马克思关于真理和认识事物的相应理论。认识事物是不断反复的过程。素描的过程就是先画事物的轮廓，然后再不断丰富细节。

\paragraph{正确做题}
做题是很重要的。自然界是没有答案给你的。平时做一道题无疑是一次小小的探索。把自己有的想法完全写出来，这就是思维的过程，这个很重要。平时做题就是解决问题，解决问题可以动用身边的一切资源。做题就是，大胆的查找资料（不是找直接的答案）。平时做题和考试做题不同，平时做题是知识的运用和方法的训练过程，考试做题是知识的运用和方法掌握的检测过程，本质是不一样的。



\newpage
\section{基本技能}
{\bf 编程语言学习：一本书在手（了解基本结构）+网上搜索资料}

\subsection{Zotero文献管理安装与简介}
现场演示！

\newpage
\subsection{\LaTeX 安装与简介}
现场演示！ Texlive 与 Texworks 或 overleaf！

\newpage
\subsection{Python 安装与简介}
Anaconda套装！现场演示！


\newpage
\section{科研与考研}
{\bf 明确定位很重要！}
\begin{itemize}
\item 拿学位找工作和继续深造拿博士学位，完全不一样！！！

\item 是否要继续博士？很关键，直接影响你自己的方向和导师培养。
\end{itemize}


\newpage
\section{自由讨论}

















\end{document}